\Subhead{1.1.1}{Линейные колебательные системы с одной степенью свободы}
Одна степень свободы означает, что положение системы задаётся одной независимой координатой $x$; состояние осциллятора задаёт пара $(x,\dot x)$. Вблизи минимума потенциальной энергии возвращающая сила линейна. Модель с вязкими потерями:
\begin{equation}\label{eq:linear}
 m\ddot x+2m\gamma\dot x+kx=F(t),\qquad \omega_0^2=k/m.
\end{equation}
Здесь $m$ --- инерция, $k>0$ --- жёсткость, $\gamma>0$ --- декремент амплитуды. Для последовательного RLC-контура $x\to q$, $m\to L$, $2m\gamma\to R$, $k\to C^{-1}$, $F\to\mathcal E$. Суперпозиция решений обусловлена линейностью, а не тем, что колебания обязательно синусоидальны.

Без силы при $\gamma<\omega_0$ получаем $x=A_0\e^{-\gamma t}\cos(\omega_dt+\phi_0)$, $\omega_d=\sqrt{\omega_0^2-\gamma^2}$. При критическом и большем затухании возврат апериодический. Без потерь фазовые траектории --- эллипсы постоянной энергии; с потерями они стягиваются к равновесию.
\begin{equation}
 E=\frac{m\dot x^2+kx^2}{2},\quad
 \dot E=F\dot x-2m\gamma\dot x^2,\quad
 Q=2\pi\frac{E}{\Delta E_{\rm пот}}\simeq\frac{\omega_0}{2\gamma}.
\end{equation}
Последнее равенство относится к слабому затуханию: за период энергия меняется мало и убывает вдвое быстрее амплитуды. Большая добротность означает малые потери, узкий резонанс и медленное установление.
